\documentclass[11pt]{article}
\usepackage{amsmath, amssymb, amsthm}
\usepackage{geometry}
\geometry{margin=1in}

\newtheorem{theorem}{Theorem}
\newtheorem{lemma}[theorem]{Lemma}
\newtheorem{corollary}[theorem]{Corollary}
\theoremstyle{definition}
\newtheorem{definition}[theorem]{Definition}
\newtheorem{remark}[theorem]{Remark}
\newtheorem{example}[theorem]{Example}

\title{The First $q$-Shifted Pair Theorem:\\
$\det M_R^{(3)} = (\det M_R^{(2)})^2$ in the Hecke Staircase}
\author{Clio}
\date{2026-04-18}

\begin{document}
\maketitle

\begin{abstract}
Let $\Pi_q = B_1 B_2 \cdots B_{n-1}$ be the staircase element in the Iwahori--Hecke algebra
$H_q(S_n)$, where $B_j = (1+T_j)(1+T_{j-1}) \cdots (1+T_1)$; set
$\Pi_q^{(k)} = B_1 \cdots B_k$ for $k \leq n-1$, and let
$M_R^{(k)}[u,v]$ be the coefficient of $T_u$ in $T_v \cdot \Pi_q^{(k)}$,
restricted to the descent-paired set $D_n \subset S_n$.
We prove that for every $n \geq 4$,
\[
\det M_R^{(3)}[D_n, D_n] \;=\; \bigl(\det M_R^{(2)}[D_n, D_n]\bigr)^2
\]
as polynomials in $\mathbb{Z}[q]$. As a corollary, the paired-ratio
$R_3(n,q)/R_2(n,q) = q^{|D_n|}$ where $|D_n| = n!/2^{\lfloor n/2 \rfloor}$
and $R_k := \det M_R^{(k)}/\det M_R^{(k-1)}$. The proof is by a
\emph{locality-reduction} argument: every $\tau_3$-block of $D_n$ is
canonically isomorphic to $D_4$ under an order-preserving identification,
reducing the general identity to a direct $6 \times 6$ computation at $n=4$.
\end{abstract}

\section{Setup}

\begin{definition}[Hecke conventions]
Let $H_q(S_n)$ be the Iwahori--Hecke algebra with generators $T_1,\ldots,T_{n-1}$
and relations $T_i^2 = (q-1)T_i + q$ and the usual braid relations.
Basis: $\{T_w : w \in S_n\}$. Right multiplication by $T_i$:
\[
T_w \cdot T_i = \begin{cases} T_{ws_i}, & w(i) < w(i+1),\\
q\, T_{ws_i} + (q-1) T_w, & w(i) > w(i+1), \end{cases}
\]
where $ws_i$ swaps positions $i$ and $i+1$ in the one-line notation of $w$.
\end{definition}

\begin{definition}[Staircase and matrices]
$B_j := (1+T_j)(1+T_{j-1}) \cdots (1+T_1)$ and $\Pi^{(k)}_q := B_1 B_2 \cdots B_k$.
The full staircase is $\Pi_q := \Pi^{(n-1)}_q$.
\end{definition}

\begin{definition}[Descent-paired set]
$D_n = \{w \in S_n : w(2i-1) > w(2i) \text{ for all } i = 1, \ldots, \lfloor n/2 \rfloor\}$.
One has $|D_n| = n!/2^{\lfloor n/2 \rfloor}$.
\end{definition}

\begin{definition}[Staircase matrix]
For $u, v \in D_n$ and $k \in \{1, \ldots, n-1\}$, set
\[
M^{(k)}[u, v] \;:=\; [T_u]\bigl(T_v \cdot \Pi^{(k)}_q\bigr),
\]
the coefficient of $T_u$ in the expansion of $T_v \cdot \Pi^{(k)}_q$.
This defines a $|D_n| \times |D_n|$ matrix over $\mathbb{Z}[q]$.
\end{definition}

We write $\tau_k(w) := (w(k+2), \ldots, w(n))$ for the length-$(n-k-1)$ suffix.
The $\tau_3$-blocks of $D_n$ are the preimages of $\tau_3$ on $D_n$; we show
below that $\Pi^{(k)}_q$ preserves this partition for $k \leq 3$.

\section{Main Theorem}

\begin{theorem}[First $q$-Shifted Pair]
\label{thm:main}
For every $n \geq 4$,
\[
\det M^{(3)}[D_n, D_n] \;=\; \bigl(\det M^{(2)}[D_n, D_n]\bigr)^2 \quad \text{in } \mathbb{Z}[q].
\]
\end{theorem}

\begin{corollary}
\label{cor:ratio}
With $R_k(n, q) := \det M^{(k)}[D_n, D_n] / \det M^{(k-1)}[D_n, D_n]$,
\[
R_3(n, q) \;=\; q^{|D_n|} \cdot R_2(n, q), \qquad |D_n| = \frac{n!}{2^{\lfloor n/2 \rfloor}}.
\]
\end{corollary}

\begin{proof}[Proof of Corollary \ref{cor:ratio} given Theorem \ref{thm:main}]
$R_3/R_2 = \det M^{(3)} \det M^{(1)} / (\det M^{(2)})^2 = \det M^{(1)}$ by Theorem~\ref{thm:main}.
For $w \in D_n$, $w(1) > w(2)$, so $T_w \cdot (1+T_1) = T_w + q T_{ws_1} + (q-1) T_w = q T_w + q T_{ws_1}$;
the $T_w$-coefficient is $q$. Since $ws_1 \not\in D_n$ (as $ws_1(1) = w(2) < w(1) = ws_1(2)$),
no off-diagonal contributions land back on $D_n$; therefore $M^{(1)}[D_n, D_n]$ is diagonal
with entries all equal to $q$, and $\det M^{(1)} = q^{|D_n|}$.
\end{proof}

\section{Proof of Theorem \ref{thm:main}}

The proof consists of three locality lemmas plus a base case at $n=4$.

\subsection{Locality}

\begin{lemma}[Locality of $\Pi^{(k)}$ for $k \leq 3$]
\label{lem:loc1}
For $k \leq 3$, $\Pi^{(k)}_q$ lies in the subalgebra $\langle T_1, T_2, T_3 \rangle \subset H_q(S_n)$.
In particular, for any $w \in S_n$, the support of $T_w \cdot \Pi^{(k)}_q$ in the $T$-basis
consists of $T_{w'}$ with $w'(i) = w(i)$ for all $i \geq 5$.
\end{lemma}

\begin{proof}
By definition $\Pi^{(k)}_q = B_1 \cdots B_k$ where $B_j$ is a product of factors
$(1+T_i)$ with $i \leq j \leq k \leq 3$, so all factors use only $T_1, T_2, T_3$.
The right-action of $T_i$ on $T_w$ produces terms supported on $\{T_w, T_{ws_i}\}$;
when $i \leq 3$, $ws_i$ swaps positions $i, i+1$ in one-line notation, leaving positions
$5, \ldots, n$ unchanged. Iterating through the factors of $\Pi^{(k)}_q$ preserves this property.
\end{proof}

\begin{lemma}[$\tau_3$-block structure]
\label{lem:blocks}
For each descent-paired suffix $\sigma' = (a_5, \ldots, a_n)$ (i.e., satisfying
$a_{2i-1} > a_{2i}$ for all $i \geq 3$), the $\tau_3$-block
$V_{\sigma'} := \{w \in D_n : \tau_3(w) = \sigma'\}$ has exactly $6 = |D_4|$ elements,
indexed by the $|D_4|$ descent-paired arrangements of the prefix-value set
$P_{\sigma'} := \{1, \ldots, n\} \setminus \{a_5, \ldots, a_n\}$.
\end{lemma}

\begin{proof}
The defining descent-pair constraints of $D_n$ split into prefix constraints
($w(1) > w(2)$ and $w(3) > w(4)$) and suffix constraints ($a_{2i-1} > a_{2i}$ for $i \geq 3$).
For fixed valid $\sigma'$ the suffix constraints hold by hypothesis; the prefix is an
arbitrary arrangement of $P_{\sigma'}$ subject to the two prefix constraints, yielding
$4!/2^2 = 6$ choices.
\end{proof}

\begin{lemma}[Block isomorphism]
\label{lem:iso}
Let $\sigma'$ be a valid suffix and $P_{\sigma'} = \{p_1 < p_2 < p_3 < p_4\}$ its prefix-value
set. Let $\phi_{\sigma'} : P_{\sigma'} \to \{1, 2, 3, 4\}$ be the unique order-preserving
bijection. Define the linear map
\[
\Phi_{\sigma'} : \operatorname{span}(T_w : w \in V_{\sigma'}) \;\longrightarrow\; \operatorname{span}(T_u : u \in D_4),
\qquad T_w \mapsto T_{\phi_{\sigma'} \circ (w|_{\{1,2,3,4\}})}.
\]
Then $\Phi_{\sigma'}$ is a $\mathbb{Z}[q]$-module isomorphism that intertwines right-multiplication
by $\Pi^{(k)}_q$ for every $k \leq 3$.
\end{lemma}

\begin{proof}
\emph{Bijection.} By Lemma~\ref{lem:blocks}, $V_{\sigma'}$ is in bijection with the
descent-paired arrangements of $P_{\sigma'}$; $\phi_{\sigma'}$ then gives a bijection
with $D_4$ (descent-paired arrangements of $\{1,2,3,4\}$), since $\phi_{\sigma'}$ is
order-preserving.

\emph{Intertwining.} It suffices to check intertwining with right-multiplication by each
$T_i$, $i \in \{1, 2, 3\}$. For $w \in V_{\sigma'}$:
if $w(i) < w(i+1)$, then $T_w T_i = T_{ws_i}$; since $\phi_{\sigma'}$ preserves order,
$\phi_{\sigma'}(w(i)) < \phi_{\sigma'}(w(i+1))$, so on the $D_4$ side
$\Phi_{\sigma'}(T_w) \cdot T_i$ uses the same rule. Since $i \leq 3$, $ws_i$ only changes
positions $i, i+1$ (both in $\{1,2,3,4\}$), so
$\Phi_{\sigma'}(T_{ws_i}) = \Phi_{\sigma'}(T_w) \cdot T_i$ with coefficient $1$. The
$w(i) > w(i+1)$ case is analogous, with the same coefficients $q$ and $q-1$. The result
extends to products by linearity and associativity.
\end{proof}

\subsection{Reduction to $n = 4$}

\begin{lemma}[Global factorization]
\label{lem:global}
For $k \in \{1, 2, 3\}$,
\[
\det M^{(k)}[D_n, D_n] \;=\; \bigl(\det M^{(k)}[D_4, D_4]\bigr)^{N(n)},
\]
where $N(n) = |D_n|/6 = \binom{n}{4} \cdot |D_{n-4}|$ is the number of valid $\tau_3$-blocks.
\end{lemma}

\begin{proof}
Since $\Pi^{(k)}_q$ fixes positions $5, \ldots, n$ (Lemma~\ref{lem:loc1}), it preserves the
$\tau_3$-block partition of $D_n$; hence $M^{(k)}[D_n, D_n]$ is block-diagonal with respect
to $\{V_{\sigma'}\}_{\sigma'}$, giving
\[
\det M^{(k)}[D_n, D_n] \;=\; \prod_{\sigma' \text{ valid}} \det M^{(k)}\bigl|_{V_{\sigma'}}.
\]
By Lemma~\ref{lem:iso}, $\det M^{(k)}|_{V_{\sigma'}} = \det M^{(k)}[D_4, D_4]$ for every
valid $\sigma'$; the count of valid suffixes equals $N(n)$ (a descent-paired sequence of
length $n - 4$ using any $(n-4)$-subset of $\{1,\ldots,n\}$).
\end{proof}

\subsection{Base case: $n = 4$}

\begin{lemma}
\label{lem:base}
$\det M^{(3)}[D_4, D_4] = \bigl(\det M^{(2)}[D_4, D_4]\bigr)^2 = q^{18}(q+1)^6(2q+1)^2$.
\end{lemma}

\begin{proof}
Direct computation (exact SymPy over $\mathbb{Q}(q)$; see
\texttt{/home/clio/projects/scratch/xmx\_final.py}).
The $\tau_2$-block decomposition of $D_4$ has sizes $(3, 2, 1)$ (indexed by $w(4) = 1, 2, 3$);
by Rule~B, $M^{(2)}[D_4, D_4]$ is block-diagonal with determinants
\begin{align*}
\det M^{(2)}|_{\text{size 3}} &= q^5(q+1)^2, \\
\det M^{(2)}|_{\text{size 2}} &= q^3(2q+1), \\
\det M^{(2)}|_{\text{size 1}} &= q(q+1).
\end{align*}
Their product is $\det M^{(2)}[D_4, D_4] = q^9 (q+1)^3 (2q+1)$.

The single $\tau_3$-block of $D_4$ is all of $D_4$ (size $6$); direct computation of the
$6 \times 6$ determinant yields
\[
\det M^{(3)}[D_4, D_4] = q^{18}(q+1)^6(2q+1)^2 = \bigl(q^9(q+1)^3(2q+1)\bigr)^2.
\]
\end{proof}

\subsection{Assembly}

\begin{proof}[Proof of Theorem \ref{thm:main}]
By Lemma~\ref{lem:global} applied to $k=2$ and $k=3$:
\begin{align*}
\det M^{(3)}[D_n, D_n] &= \bigl(\det M^{(3)}[D_4, D_4]\bigr)^{N(n)}, \\
\det M^{(2)}[D_n, D_n] &= \bigl(\det M^{(2)}[D_4, D_4]\bigr)^{N(n)}.
\end{align*}
By Lemma~\ref{lem:base}, $\det M^{(3)}[D_4, D_4] = (\det M^{(2)}[D_4, D_4])^2$. Raising both
sides to the $N(n)$-th power gives
\[
\det M^{(3)}[D_n, D_n] = \bigl(\det M^{(2)}[D_n, D_n]\bigr)^2.
\]
\end{proof}

\section{Computational verification}

All steps of the proof were verified by exact symbolic computation over
$\mathbb{Z}[q]$ in SymPy:
\begin{itemize}
\item $n = 4$ (Lemma~\ref{lem:base}): direct $6 \times 6$ determinant.
See \texttt{xmx\_final.py}; output matches $\det M^{(3)} = q^{18}(q+1)^6(2q+1)^2$.

\item $n = 5$ (locality): all $5$ valid $\tau_3$-blocks of $D_5$ are IDENTICAL
as $6 \times 6$ matrices after the order-preserving relabeling of $\phi_{\sigma'}$.
Each has determinant $q^{18}(q+1)^6(2q+1)^2$, and the collective $\det M^{(2)}$
equals $(q^9(q+1)^3(2q+1))^5$.

\item $n = 5, 6, 7$ (main identity): $\det M^{(3)}[D_n] - (\det M^{(2)}[D_n])^2 = 0$
as polynomials (verified by \texttt{sp.expand} and direct comparison).
See \texttt{prove\_locality\_check.py}.
\end{itemize}

\section{Remarks}

\begin{remark}[Correction of the original conjecture]
The originally conjectured form
$R_3(n, q) = q^{6\binom{n}{4}} R_2(n, q)$
is numerically consistent with Corollary~\ref{cor:ratio} for $n \in \{4, 5, 6\}$
because $|D_n| = 6 \binom{n}{4}$ for these values, but the two expressions
\emph{diverge} at $n = 7$ and beyond:
\[
|D_7| = 630, \qquad 6\binom{7}{4} = 210.
\]
Explicit computation confirms $R_3(7,q)/R_2(7,q) = q^{630}$, not $q^{210}$.
The combinatorial cleanliness of the statement $\det M^{(3)} = (\det M^{(2)})^2$
(and its corollary with exponent $|D_n|$) appears to be the natural formulation.
\end{remark}

\begin{remark}[Why the identity $\det M^{(3)}[D_4] = (\det M^{(2)}[D_4])^2$ at $n=4$?]
Theorem~\ref{thm:main} reduces to a direct computational check at $n=4$.
The numerological coincidence that $\Pi^{(3)}_q$ for $n=4$ equals the full staircase
$\Pi_q$ of $H_q(S_4)$, while $\Pi^{(2)}_q$ is the one-step-short truncation, suggests
a deeper structural reason --- perhaps related to the eigenvalue palindromy of the
staircase operator, or to a duality between $\Pi^{(k)}$ and $\Pi^{(n-1-k)}$. We leave
this as an open question (Section~\ref{sec:open}).
\end{remark}

\begin{remark}[Higher $k$]
The locality argument extends: for any fixed $k$, $\Pi^{(k)}_q$ lies in
$\langle T_1, \ldots, T_k \rangle \cong H_q(S_{k+1})$, and the $\tau_k$-block structure
of $D_n$ reduces to $D_{k+1}$ under order-preserving relabeling. Hence all identities
among $\{\det M^{(j)}[D_n]\}_{j \leq k}$ reduce to identities at $n = k+1$. This suggests
a systematic program for higher pairs $R_{2j}/R_{2j+1}$.
\end{remark}

\section{Open questions}
\label{sec:open}

\begin{enumerate}
\item \textbf{Structural reason for the $n=4$ squaring.} Is there a proof of
$\det M^{(3)}[D_4] = (\det M^{(2)}[D_4])^2$ that avoids direct computation?
Eigenvalue palindromy and the bicoboundary (Apr 17) result may be relevant.
\item \textbf{Higher pairs.} What is the analogous identity for $\det M^{(5)}$ vs.\
$\det M^{(4)}$? By the locality program it reduces to $n = 6$, but $|D_6| = 90$;
exact computation is feasible but nontrivial.
\item \textbf{Non-paired ratios.} What is $R_2/R_3$ as a polynomial in $q$ for
generic $n$? Our corollary gives $R_3 = q^{|D_n|} R_2$, but the individual $R_k$'s
have rich structure involving $(q+1)$ and $(2q+1)$ factors with multiplicities
$3 |D_n|/6$ and $|D_n|/6$ respectively.
\end{enumerate}

\end{document}
